\title{X-LoRA: Mixture of Low-Rank Adapter Experts, a Flexible Framework for Large Language Models with Applications in Protein Mechanics and Molecular Design}

\begin{abstract}
We present a mixture of expert approach for fine-tuning large language models using a deep, layer-wise, token-level methodology based on low-rank adaptation (LoRA). Beginning with a collection of pre-trained LoRA adapters, our gating mechanism leverages hidden states to dynamically combine adapted layers, enabling the resulting X-LoRA model to utilize diverse capabilities and form novel deep layer-wise combinations for task-solving. This design draws inspiration from biological principles of universality and diversity, where neural network components are reused in multiple hierarchical forms. Consequently, the X-LoRA model can be easily integrated with any existing large language model (LLM) without requiring alterations to its core architecture. We have developed a specialized X-LoRA model tailored to scientific applications, including forward and inverse analysis tasks and improved reasoning abilities, with a focus on biomaterial analysis, protein mechanics, and design. The significance of this work lies in providing accessible, extensible, and adaptable models equipped with strong domain expertise and the ability to synthesize knowledge across fields. Incorporating experts in biology, mathematics, reasoning, bio-inspired materials, mechanics, chemistry, protein biophysics, and quantum-mechanics-based molecular properties, we perform a suite of physics-centered case studies. These include knowledge recall, protein mechanics forward/inverse challenges, protein design, adversarial agentic modeling encompassing ontological knowledge graph creation, and molecular design. The model is capable not only of quantitatively predicting nanomechanical properties of proteins or quantum mechanical molecular characteristics but also of reasoning about these outcomes and accurately inferring plausible mechanisms that explain distinct molecular behaviors.
\end{abstract}

\section{Introduction}  
Large language models (LLMs)~\cite{Vaswani2017AttentionNeedc,Touvron2023LlamaModels,OpenAI2023GPT-4Report,Chowdhery2022PaLM:Pathways,Jiang2023Mistral7B,Gunasekar2023TextbooksNeed} have become highly popular, including for creating specialized models that excel in specific task types, reasoning processes, or scientific fields~\cite{Bubeck2023SparksGPT-4,Buehler2023MechGPTModalities,Nejjar2023LLMsAnalysis,Buehler2023GenerativeDesign,Luu2023BioinspiredLLM:Materials,Luu2023GenerativeSolvents,Buehler2023MeLMProblemsb,Ge2023OpenAGI:Experts}. However, the training of these models can be expensive, particularly when a wide variety of capabilities are required. Techniques such as low-rank adapters (LoRA)\cite{hu2021lora} have been introduced as a more resource-efficient alternative, but these adaptations tend to focus on narrower domains of knowledge. The core idea behind LoRA is to use low-rank matrices that augment the original full-scale matrix, with only these low-rank matrices being trainable components of the model. Because only the adapter layers are trainable, these models generally retain the information learned during the pre-training phase of the base model, while becoming more specialized for specific tasks and remaining computationally efficient. The efficiency of training LoRA adapters enables the creation of numerous specialized models using this approach.

Although LoRA adapters offer a resource-conscious approach to building specialized model functionalities, the question of how to combine multiple such adapters into enhanced models that provide a unified and even superior set of capabilities remains unresolved. Indeed, combining LLMs has emerged as an intriguing research area~\cite{Kim2023SOLARUp-Scaling}. Experimental approaches, including using spherically interpolated model parameters (SLERP) and other techniques~\cite{Arcee-ai/mergekit:Models.}, have demonstrated promising outcomes and the potential of model mixing to unlock new abilities. Mixed models developed through these methods are typically constructed {\it a priori} following specific algorithms and, while they have achieved encouraging benchmark performance, further investigation is needed into the systematic development of mixed models and the mechanisms by which certain features arise.

We propose that dynamically mixing parameters could serve as a more versatile approach, as it enables the combined model to investigate novel parameter combinations during inference.

Similar methodologies have employed mixture of experts (MoE) frameworks~\cite{Jacobs1991AdaptiveExperts,Hampshire1992TheRecognition,Jordan1994HierarchicalAlgorithm}, including the recently introduced Mixtral LLM~\cite{Jiang2024MixtralExperts}. Nevertheless, these methods can be computationally intensive, requiring substantial resources even at inference time. In this work, we introduce a computationally efficient mixture-of-experts technique utilizing a collection of distinct LoRA adapters that can be trained independently with ease. Our method delivers a highly granular and flexible solution, drawing inspiration from a biological model that repurposes neural network components to build a hierarchy of interrelated models, spanning from the foundational model to agentic systems (Fig.~\ref{fig:hierarchical_concept}). To rigorously evaluate its effectiveness, our study targets the application of this strategy to science-oriented LLMs, specifically addressing protein mechanics challenges within the broader context of bio-inspired materials analysis and design.

We develop an algorithm that dynamically combines adapters at the token-level state, with fine granularity allowing mixing at the individual layer level. Termed X-LoRA, this approach provides access to a broad range of capabilities vital for scientific applications of multimodal language models~\cite{Hu2023DeepScience,Buehler2023AMetamaterials,Buehler2022ProteinNetworks,}. 

\subsection{Fundamental concepts of X-LoRA}

We begin with a concise review of the principles underlying LoRA adapters. The core idea behind low-rank adaptation (LoRA)~\cite{hu2021lora} is the assumption that weight updates possess a low “intrinsic dimension,” leveraging this by freezing the original weights $W_0 \in \mathbb{R}^{d \times k}$ and limiting updates to a low-rank decomposition

\begin{equation}
W_0 + \Delta W_0 = W_0 + BA 
\end{equation}

where $B \in \mathbb{R}^{d \times r}$ and $A \in \mathbb{R}^{r \times d}$ with the rank $r \ll {\rm min}(d,k)$. In the training process with LoRA~\cite{hu2021lora}, the pretrained weights $W_0$ remain fixed while only the matrices $A$ and $B$ contain trainable parameters. The forward pass of a LoRA layer can be described by the equation:

\begin{equation}
    h = W_0x + \Delta W_0x = W_0x + BAx
\end{equation}

It is worth noting that in practice, LoRA adapters use a fixed scalar weight $\alpha$, which is applied to the decomposition matrices, resulting in:

\begin{equation}
    h = W_0x + BAx \cdot \alpha
\end{equation}

Building upon this, X-LoRA introduces the concept of employing a collection of LoRA adapters, each possessing distinct capabilities, as demonstrated in previous science-oriented LLM research (e.g., \cite{Buehler2023GenerativeDesign}). Instead of statically mixing or integrating models, we develop a dynamic gating method that scales each LoRA adapter on a per-token and per-layer basis, enabling deeper mixing within the model. More specifically, for a given LoRA adapter $i$, the original scaling parameter $\alpha_i$ is multiplied by a scaling factor $\lambda_i$, producing an adjusted scaling value $\alpha^*_i$:

\begin{equation}
    \alpha^*_i = \alpha_i \cdot \lambda_i
\end{equation}

The scaling factor $\lambda_i$ is predicted by an X-LoRA scaling head that leverages the model’s hidden states, representing the trainable part of the X-LoRA model. Because the scaling head can utilize sophisticated encodings, it can be efficiently implemented as a straightforward feed-forward neural network with one or a few layers (in the X-LoRA implementation, users can specify the architecture details of this neural network).

The scaling computation is performed individually for each token in the sequence, providing fine-grained control as every LoRA adapter active at specific layers of the LLM is scaled accordingly.

From this point onward, we refer to this method of gating the components of the adapters in this manner as scaling. Additional information about this method is given in the Materials and Methods section.

\subsection{Paper outline}

The structure of this paper is organized as follows. Initially, we outline the methodology and training strategy, including the pertinent datasets, which leads to the creation of an X-LoRA model with specialized capabilities in the physical sciences, particularly emphasizing biomaterials such as proteins. Subsequently, we demonstrate a series of experiments applying the X-LoRA to diverse tasks, including question answering, conversational and agentic modeling, protein design and analysis, among others. A thorough examination of the scaling behaviors is carried out, and the approach is validated by comparing it with molecular modeling and various other physical data and methodologies.

\section{Results and discussion}

The development of the X-LoRA model follows these steps:

\begin{enumerate}
    \item Training the foundational base LLM (completed in prior work) \item Separately training a collection of adapters to acquire expertise across several domains (these represent distinct or overlapping skill sets and knowledge areas) \item Training the integrated X-LoRA model
\end{enumerate}

Our experimental work begins with the training of multiple LoRA adapters. We create nine adapters, each fine-tuned for specific expertise, based on the Zephyr-7B-$\beta$ model~\cite{HuggingFaceH4/zephyr-7b-betaFace}, which itself is built upon the Mistral-7B model\cite{Jiang2023Mistral7B}:

\begin{enumerate}
    \item Bioinspired materials \item Chain-of-thought (CoT) and reasoning \item Chemistry \item Mathematics \item Physics \item Biology \item Mechanics and materials \item Platypus: Logic and reasoning \item Protein mechanics tasks (generative sequence-to-property and inverse capabilities)
\end{enumerate}

These specific domains of focus and expertise were selected because our ultimate goal is to develop an LLM centered on STEM, with an emphasis on materials science. Nonetheless, we hypothesize that our methodology can be broadly applied across various fields, particularly to generate modalities that foster links between different areas of knowledge~\cite{Buehler2023MechGPTModalities,Buehler2023MeLMProblemsb}.

Following the creation of the LoRA adapters, we proceed to train the X-LoRA model. This involves utilizing a subset of the combined training samples originally used for the individual LoRA adapters (for more details, see Materials and Methods, which includes comprehensive descriptions of the dataset, implementation, and related aspects).

The X-LoRA system serves as an experimental platform to evaluate its effectiveness in multiple scenarios. We perform a series of tests, such as tackling complex question answering, applying the generative model to protein design, conducting analysis and reasoning, and developing more advanced interaction strategies via adversarial agentic modeling. Throughout these analyses, we compare the performance of the base model against the X-LoRA model, while monitoring how the scaling head leverages various layers of the different experts during the generation or formulation of responses.

\subsection{Question answering and observed X-LoRA layer-wise scaling weights}

Our initial investigation focuses on the distinctions between the X-LoRA model and the foundational base model. We pose the following question to both models:

\begin{quote}
\small\texttt{What material behavior (stiffening vs softening) causes supersonic crack speeds in mode II cracks? Think step by step.}
\end{quote}

The base model replies:

\begin{quote}
\small\texttt{Supersonic crack speeds, also known as Rayleigh waves, occur in mode II cracks due to the stiffening material behavior of the material near the crack tip. In mode II fracture, the loading is shear-dominated and the crack propagates perpendicular to the direction of applied stress. Here's a step-by-step explanation:

As the crack tip nears the material, the stresses surrounding it escalate rapidly. This creates a zone with elevated stress concentrations, known as the stress singularity.

Within this stress singularity, the material undergoes elastic deformation that causes an increase in stiffness, exhibiting stiffening behavior. This stiffening arises because the material adjacent to the crack tip becomes increasingly constrained as the crack advances [...]}
\end{quote}

Drawing on our own domain expertise~\cite{Buehler2008AtomisticFailure}, this response is not only inaccurate but also unnecessarily verbose. Conversely, the X-LoRA model replies:

\begin{quote}
\small\texttt{The material behavior that causes supersonic crack speeds in mode II cracks is stiffening.}
\end{quote}

This answer is accurate, straightforward, and concise~\cite{Buehler2008AtomisticFailure}. 

The contrast between the base model and the X-LoRA model is even more pronounced in specialized protein analysis and design tasks. The given task is:

\begin{quote}
\small\texttt{CalculateForceHistory< M S K S N T Y R M L V L L E D D T K I N K E D E K F L K G K P G K M H E F V D E L I L P F N V D E L D E L N T W F D K F D A E I C I P Y V E I L K E S G M K > }
\end{quote}

The base foundation model's reply is:

\begin{quote}
\small\texttt{I'm sorry but the input provided is not a valid mathematical expression. It appears to be a sequence of characters that do not form a calculable function. Please provide a valid mathematical expression for me to calculate.

However, I can tell you that the sequence of characters you provided is actually a series of function and variable names that may be part of a larger program or script in a specific programming language. Without context, it's impossible to determine the intended calculation.

If you provide more information, such as the programming language and the context in which this expression appears, I may be able to provide more insight into its meaning and possible calculations. }
\end{quote}

The response from the X-LoRA model is as follows, correctly fulfilling the design task by producing the force-deformation curve related to this protein:

\begin{quote}
\small\texttt{[0.013,0.066,0.078,0.045,0.072,0.126,0.127,0.107,0.119, 
0.074,0.106,0.098,0.103,0.128,
0.077, 0.069,0.131,0.108,0.144,0.174, 0.166,0.183,0.185,0.217,0.175,0.230,0.187,0.224, 0.267,0.301,0.315]}
\end{quote}

Figure~\ref{fig:QA_weights_sample_0} displays outcomes from this question answering along with the observed X-LoRA scaling weights. The two queries, one related to mechanics and materials (Fig.~\ref{fig:QA_weights_sample_0}a) and the other to protein mechanics (Fig.~\ref{fig:QA_weights_sample_0}b), reveal distinct activation patterns. Notably, we observe that each case employs a complex pattern of scaling values, implying that the X-LoRA model leverages a heterogeneous combination of different adapters distributed across layers. Additionally, the heatmap in Fig.~\ref{fig:QA_weights_sample_0}a exhibits a wider range of activation spanning many experts and layers, including certain regions of strong activation. In contrast, the heatmap in Fig.~\ref{fig:QA_weights_sample_0}b demonstrates a more focused activation pattern, with fewer widespread high values. We interpret this as possibly indicating that the gating function on the right activates experts in a more selective manner.

In Fig.~\ref{fig:QA_weights_sample_0}, the areas that appear as bright lines indicate paths of high activation, suggesting that the scaling function is emphasizing certain experts more strongly for particular layers. The observed patterns indicate that the importance is not uniformly spread across all experts and layers; rather, it tends to be sparse and selective. This sparsity aligns with findings from other mixtures-of-experts models, where only a subset of experts is activated for each input to specialize in different segments of the problem space.

When comparing the maps across various tasks, both the sparsity and the distinct patterns may be crucial for gaining deeper insights into which experts the model depends on for different input types or at different stages of processing within the neural network (we will later discuss how these maps evolve dynamically as complex tasks are solved).

To further investigate the model's behavior, we examine a broader range of questions to understand how the X-LoRA model employs different experts given distinct prompts. Our focus is on determining whether, and how, the X-LoRA model integrates multiple adapters when questions overlap across different skills or domains. Fig.~\ref{fig:QA_weights} presents results from question answering alongside the observed X-LoRA scaling weights, corresponding to the queries listed in Table~\ref{tab:table_qa}. Table~\ref{tab:table_qa} includes responses from both X-LoRA and the baseline model, Zephyr-7B-$\beta$. The contrasts between the two models are notable.

To highlight this visually, incorrect or questionable responses are marked with a \redhighlight{red highlight}. The comparison clearly demonstrates that X-LoRA delivers much higher accuracy and more concise answers.

As anticipated, we detect a complex blending of adapters, frequently involving the activation of multiple dominant LoRA experts. For example, Fig.~\ref{fig:QA_weights}d shows that a question about the mechanics of pine cones triggers strong activation of the mechanics/materials adapter combined with the CoT/reasoning and bioinspired adapters. Although X-LoRA selects the correct letter, the answer is not entirely accurate since pine cones open in dry, rather than wet, conditions. However, a subsequent question probing whether release occurs in dry or wet conditions yields the correct response.

In a protein-related query, Fig.~\ref{fig:QA_weights}f reveals adapter mixing in reply to a question concerning the mechanics of Bombyx mori silk, activating a combination of bioinspired, mechanics/materials, and biology adapters. For highly specialized tasks like protein design (Fig.~\ref{fig:QA_weights}a), we observe predominantly singular activation of one expert.

A more detailed examination of the heatmap illustrating scalings in Fig.~\ref{fig:QA_weights}a reveals that beyond the clearly prominent importance of the protein mechanics expert, there are distinct bands where certain experts receive significant emphasis across multiple layers. This pattern is especially evident for the mechanics/materials and reasoning experts. These bands appear intermittently and display a fluctuating pattern, which may indicate particular conditions or features that these experts process efficiently. Notably, the reasoning expert (positioned immediately to the left of the protein mechanics expert) is prominent, showing several layers with the highest activation, particularly in the upper layers (approximately layers 100 to 140). We hypothesize that this may suggest a critical function in the model’s ultimate decision-making or output generation, though further investigation is required to confirm this. Additionally, there is a more localized yet strong activation for the physics and biology experts at specific layers, potentially reflecting these experts’ specialization in features or representations associated with those layers.

Moreover, Fig.~\ref{fig:QA_token_history} presents a token-level history aggregated over all layers’ scalings. This analysis highlights the histories of two examples: Fig.~\ref{fig:QA_token_history}a relates to a task involving pine cone seed release, and Fig.~\ref{fig:QA_token_history}b corresponds to a sequence of protein mechanics tasks. The protein mechanics task involved multiple prompt-answer interactions: initially two property calculations followed by a generative task aimed at designing a new protein towards the conclusion. Although the use of experts remains generally consistent, noticeable shifts occur in the utilization of the LoRA experts throughout the process. In particular, the protein mechanics expert becomes highly concentrated during the final phases of generation when a new protein is being created.

Table~\ref{tab:table_qa} illustrates various use cases across multiple domains and provides a comparison to the performance of the base model alone. We delve into several of these examples in greater detail.

For example, we present a question that combines elements of protein science and biology with logical reasoning, specifically concerning protein expression patterns. X-LoRA provides the correct answer. In contrast, the base model’s response contains several inaccuracies and logical errors when evaluating the marker protein expressions for each cell type. The base model accurately determines that fibroblasts express Protein B and do not express Protein A, based on the mutual exclusivity between Protein A and Protein B. However, it fails to explicitly state that fibroblasts also express Protein C. Since there is no rule preventing Protein C expression alongside Protein B, and fibroblasts have no restriction against Protein C, it is logical that fibroblasts express Protein C as well. Regarding neurons, the base model correctly starts by asserting that neurons express Protein A and not Protein B due to the mutual exclusivity rule. Nevertheless, its explanation regarding whether neurons express Protein C or both Protein C and Protein B is confusing and contains an error. Given that neurons express Protein A and do not express Protein B (as specified), the only remaining possibility according to the rules is that neurons express Protein C. There is no scenario where neurons express Protein B, contrary to what was suggested in the initial analysis. Therefore, neurons express Protein A and Protein C, but not Protein B. For epithelial cells, the base model correctly identifies that they do not express Protein A but do express both Protein B and Protein C. However, the reasoning provided is convoluted and includes incorrect claims about mutual exclusivity between Protein A and Protein C, which was not part of the original premises. The straightforward explanation is that epithelial cells do not express Protein A but must express the remaining two proteins, Protein B and Protein C, because the initial rules only restrict co-expression of Protein A and Protein B.

In contrast, X-LoRA accurately identifies the marker proteins expressed by each cell type based on the given observations and logical reasoning. Regarding fibroblasts, X-LoRA correctly infers that fibroblasts do not express Protein A because the presence of Protein A excludes Protein B expression. Since fibroblasts do express Protein B, and considering that each cell type expresses a distinct combination of three marker proteins (an incorrect assumption introduced, as the original statement only indicated that each cell type expresses a unique combination without specifying the number of proteins), the conclusion that fibroblasts also express Protein C is valid based on the information provided (apart from the misunderstanding about the number of proteins). For neurons, the model accurately identifies expression of Protein A and Protein C. The reasoning is logical: Protein B cannot co-express with Protein A, and since neurons express Protein A and another marker protein that is not Protein B, they must express Protein C. Finally, for epithelial cells, X-LoRA correctly deduces that these cells express Protein B and Protein C, since epithelial cells do not express Protein A and, by elimination, the two expressed marker proteins are Protein B and Protein C.

Regarding the question of pine cone opening in humid versus dry environments, the base model provides a confusing and factually incorrect response, whereas X-LoRA delivers a clear and precise explanation. Similar examples include the question about the exceptional mechanical properties of Bombyx mori silk, where the base model erroneously claims that silk fibers conduct electricity. The base model also fails on several other specialized domain tasks; for instance, when asked about the role of hyperelasticity in maximum fracture speeds, X-LoRA gives a concise and accurate response while the base model cannot answer. In another case concerning the strong adhesion of gecko feet, X-LoRA correctly highlights the significance of setae and nanoscale effects, while the base model incorrectly attributes adhesion to special-purpose glue proteins. These and other instances presented in the table and throughout the paper provide compelling evidence for X-LoRA’s superior capabilities compared to the base model.

Figure~\ref{fig:perf_comp_bioinspiredexam} presents a comparison of knowledge recall evaluation experiments conducted using the bio-inspired knowledge test set introduced in \cite{Luu2023BioinspiredLLM:Materials}. It is evident that X-LoRA delivers the best performance, despite being a considerably smaller model compared to the BioinspiredLLM or Llama-BioLLM models (7B parameters for X-LoRA versus 13B for the others). The figure also includes comparisons with other recently released large language models, namely Orca-13B and Llama-13b-chat.

Additionally, the plot contrasts X-LoRA with its base model, the Zephyr-7B-$\beta$, highlighting that X-LoRA achieves markedly better results. These findings demonstrate that strong performance is attainable with X-LoRA even when smaller base models are employed.

Figure~\ref{fig:image_synth} illustrates an application in image synthesis, where X-LoRA is utilized to generate a prompt for other models, such as DALL-E 3, as shown here.

\begin{quote}
\small\texttt{Create a prompt that I can use for image generation that accurately describes the microstructure of a hierarchical composite.

Explicitly describe the geometry.}
\end{quote}

The comparison in Figure~\ref{fig:image_synth}(b)-(c) clearly reveals a more succinct prompt, which in turn leads to a significantly more precise image generated by DALL-E 3~\cite{DALLECard}. Notably, the prompt generated by Zephyr-7B-$\beta$ fails to accurately capture the given instruction and introduces various additional design elements, such as fibers and nanoparticles. These details were absent from the original prompt, which specifically focused on the microstructure of a hierarchical composite. The X-LoRA prompt, being more concise and targeted, yields a substantially improved outcome.

\subsection{Generative solutions to protein analysis}

We now turn to specific protein generative and analytical tasks. These capabilities emerge within the X-LoRA model due to the protein mechanics adapter and are further strengthened by integrating the model’s knowledge in biology, bio-inspired materials, reasoning, and logical deduction. This section concentrates on examining protein-related tasks independently, providing a quantitative evaluation of this functionality.

Fig.~\ref{fig:protein_force_history_prediction} presents findings from the protein task aimed at predicting force-deformation behaviors based on sequences. Fig.~\ref{fig:protein_force_energy_prediction} illustrates the overall accuracy in predicting unfolding force and unfolding energy from the sequence, comparing actual values with predictions (yielding $R_2=0.85$). In general, the model demonstrates outstanding performance and, as will be discussed in the following section, it also excels in the inverse problem of design and cycle consistent evaluation\cite{Lu2023GenerativePropertiesb}. However, unlike previous studies that relied on highly specialized GPT-type models for such tasks, our X-LoRA model integrates these task-specific capabilities with a wide range of other deep specializations achieved through various fine-tuned adaptations.

\subsection{Protein design and analysis}

We now extend the testing scenarios to more comprehensive applications that encompass multiple capability areas. Specifically, the experiments described in this section demonstrate how the X-LoRA model leverages a unique set of competencies and how agentic modeling can generate more complex and thorough responses. We start with a protein design task, where the model is tasked with creating a novel protein that satisfies a specified target force-deformation profile. These force-deformation responses have been investigated in previous work~\cite{Ni2023ForceGen:Model} and correspond to measurements of force versus deformation as the protein is stretched at both ends (representing a classical protein unfolding experiment~\cite{Lu1998UnfoldingSimulation.}).

Fig.~\ref{fig:design_protein} illustrates the application of the generative protein task. We create a protein exhibiting a specified force-deformation response (based on training data and boundary conditions from the original molecular dynamics (MD) simulations of protein pulling, as described in~\cite{Ni2023ForceGen:Model}), and subsequently evaluate the predicted sequence to verify if it achieves the intended target characteristics (Fig.~\ref{fig:design_protein}a). This cycle-consistent evaluation is crucial for verifying whether the inverse design capabilities align with the forward prediction results.

In Fig.~\ref{fig:design_protein}b, the predicted amino acid sequence, \texttt{MSKSNTYRMLVLLEDDTKINKEDEKFLKGKPGKMHEFVDELILPFNVDELDELNTWFDKFDAEICIPYVEILKESGMK}, is presented. Additionally, Fig.~\ref{fig:design_protein}c displays the folded conformation of the protein, obtained via AlphaFold 2~\cite{Jumper2021HighlyAlphaFoldb}. To understand the relationship of the designed protein with existing sequences, Fig.~\ref{fig:design_protein}d utilizes a Basic Local Alignment Search Tool (BLAST) Tree~\cite{Altschul1990BasicTool}. Notably, the designed protein shows connections to hypothetical proteins previously examined in earlier studies, although these had not yet been identified or structurally resolved.

By employing agentic modeling and incorporating multiple X-LoRA models as components within a multi-agent framework, we further investigate the predicted sequence using our X-LoRA model. This approach leverages the enhanced ability of our model to not only interpret protein mechanics but also integrate this knowledge with other domains. Using the prediction generated as described above as input, we request an in-depth analysis and suggestions for sequence variations. We begin with the following query:

\begin{quote}
\small\texttt{I have identified this amino acid sequence: MSKSNTYRMLVLLEDDTKINKEDEKFLKGKPGKMHEFVDEL ILPFNVDELDELNTWFDKFDAEICIPYVEILKESGMK.

How can I use it?

Provide logical steps, discuss specifically the engineering of the problem.}
\end{quote}

We engage multiple agentic interactions to develop the solution, applying adversarial strategies (details provided in Materials and Methods) to encourage the model to probe deeper and more intricate aspects of the concepts addressed. The summary of this interaction, as synthesized by the X-LoRA model (the complete conversation is available in the Supplementary Material, SI Conv Example 1), is as follows:

\begin{quote}
\small\texttt{The discussion centers on the prospective uses of a particular amino acid sequence within protein engineering. This sequence comprises both hydrophobic and hydrophilic residues, indicating its potential to form a stable conformation and engage with other molecules. To design a protein incorporating this sequence, computational approaches can predict its secondary and tertiary structures; however, challenges include the intricate nature of protein folding and the computational demands involved. Enhancements can be achieved by creating more efficient algorithms and integrating novel experimental findings. It is essential to validate computational models to guarantee their precision and applicability in real-world scenarios. Methods such as cross-validation, regularization, early stopping, ensemble techniques, and transfer learning are valuable to avoid overfitting and underfitting. Additionally, interpreting and explaining these models aids in guiding their practical implementation.}
\end{quote}

We also request the model to generate a summary of key points, resulting in:

\begin{quote}
\small\texttt{\begin{itemize}
    \item The amino acid sequence MSKSNTYRMLVLLEDDTKINKEDEKFLKGKPGKMHEFVDELI LPFNVDELDELNTWFDKFDAEICIPYVEILKESGMK holds promise for protein engineering due to its combination of hydrophobic and hydrophilic residues. Computational tools can forecast the secondary and tertiary structures, although limitations arise from the complexity of protein folding and computational resource requirements.
    \item Ensuring the reliability of computational models is vital for their effectiveness in practical uses. Approaches like cross-validation, regularization, early stopping, ensemble methods, and transfer learning help mitigate overfitting and underfitting issues.
    \item Understanding and elucidating models further supports their application in practical contexts.
\end{itemize}
}
\end{quote}

In conclusion, the main takeaway is:

\begin{quote}
\small\texttt{The primary takeaway is that the amino acid sequence MSKSNTYRMLVLLEDDTKINKEDEKFLKGKPGKMHEFVDELILPFNVDELDELNTWFDKFDAEICIPYVEILKESGMK holds promise for protein engineering applications owing to its combination of hydrophobic and hydrophilic residues. To design a protein featuring this sequence, computational techniques can predict its secondary and tertiary structures; however, challenges include the complexity inherent in protein folding and the computational resources required. Ensuring the accuracy and practical relevance of computational models is essential, which can be achieved through validation strategies such as cross-validation, regularization, early stopping, ensemble approaches, and transfer learning to avoid overfitting and underfitting. Additionally, interpreting and explaining these models can aid in guiding real-world applications.}
\end{quote}

We can further delve into detailed sequence engineering processes by prompting the model for targeted design recommendations aimed at improving protein stability:

\begin{quote}
\small\texttt{What modifications could be made to the following protein sequence to enhance its stability: MSKSNTYRMLVLLEDDTKINKEDEKFLKGKPGKMHEFVDELILPFNVDELDELNTWFDKFDAEICIPYVEILKESGMK?}
\end{quote}

A summary of the interaction is provided (full dialogue available in Supplementary Material, SI Conv Example 2):

\begin{quote}
\small\texttt{The discussion centers on strategies to increase protein stability by incorporating hydrophobic amino acids and forming disulfide bonds. The biologist explains that hydrophobic residues typically contribute to stability, and disulfide bonds can further reinforce structural integrity. The engineer requests concrete examples of hydrophobic amino acids and guidance on the appropriate number of disulfide bonds to add. The biologist cites examples of proteins successfully stabilized using these methods, such as insulin, bovine pancreatic trypsin inhibitor, and green fluorescent protein. The conversation also addresses potential challenges or downsides of adding hydrophobic amino acids or disulfide bonds, including alterations in protein function, added complexity, and risk of aggregation. The biologist recommends mitigating these issues through prudent amino acid selection, computational modeling, and empirical testing.}
\end{quote}

The list of main points:

\begin{quote}
\small\texttt{\begin{itemize}
\item To enhance the stability of a protein sequence, one can introduce hydrophobic amino acids and disulfide bonds.

Hydrophobic residues typically contribute to protein stability, and disulfide bonds also serve to strengthen the structure.

Examples of proteins successfully stabilized through these modifications include insulin, bovine pancreatic trypsin inhibitor, and green fluorescent protein.
\item Possible drawbacks or challenges of adding hydrophobic amino acids or disulfide bonds include alterations in protein function, increased structural complexity, and the risk of aggregation.
\item These challenges can be mitigated by careful amino acid selection, computational modeling, and empirical validation.
\end{itemize}
}
\end{quote}

Ultimately, the core takeaway is:

\begin{quote}
\small\texttt{The single most critical insight from the discussion is that introducing hydrophobic amino acids and disulfide bonds into the protein sequence MSKSNTYRMLVLLEDDTKINKEDEKFLKGKPGKMHEFVDELILPFNVDELDELNTWFDKFDAEICIPYVEILKESGMK will improve its stability. This directly answers the initial query.}
\end{quote}

From the viewpoint of protein engineering, these recommendations are logical as they are expected to result in more stable protein designs. Concrete examples and proposed strategies are presented for researchers to further investigate.

In a separate example, we examine three sequences. First, we request the model to evaluate the stability of each, then to reason about the results: Identify the strongest candidate and justify why that protein is likely the most stable. Additionally, we ask the model to describe the steps needed for its laboratory synthesis. The three sequences (manually derived by altering one sequence from the test set) are:

\begin{quote}
\small\texttt{
\begin{itemize}
    \item LNTWFDKFDAEICIPYVEILKESGMK \item AITWFDKFDAEICIPYVEALKIAAMK \item AAAAAAAAKFDAAEICIIAAAAAAAAKIAAMK
\end{itemize}
}
\end{quote}

This example assesses whether the model can integrate various abilities and adapt dynamically throughout the conversation. Specifically, it tests the capacity to switch fluidly between protein mechanics, rational analysis, hypothesis generation, and related functions.

Figure~\ref{fig:conversation_evolve} summarizes the conversation, illustrating how the scalings evolve throughout the discussion (the analysis includes both the deep layer-wise adapter scaling and an aggregated assessment to quantify the effective contribution of the different experts). The findings clearly indicate that the protein task is initially highly significant, progressively giving way to a stronger influence from the bioinspired adapter and the biology adapter. By the conclusion of the conversation, the biology adapter emerges as the dominant contributor.

To evaluate the predictions and assess the quality of the responses, the three proteins examined in this task are shown in Figure~\ref{fig:protein_visuals} (structures folded using Alpha Fold 2~\cite{Jumper2021HighlyAlphaFoldb} through ColabFold \cite{Mirdita2022ColabFold:All}). We observe that the protein with the most stable structure in the center exhibits the most well-defined secondary structure, consistent with the idea that it possesses the highest unfolding force, as predicted by the X-LoRA model in Figure~\ref{fig:conversation_evolve}. The other two proteins display less organization; the first is predominantly helical with some turns, while the last one is partially unstructured and contains a geometric kink. These structural characteristics account for the lower unfolding forces predicted by the model.

By examining the visual depictions of the folded structures, we verify that the most stable configuration in the center exhibits the highest degree of organized secondary structure, consistent with the idea that it produces the greatest unfolding force as predicted by the model. We observe that the other two configurations are less structured. The first is predominantly helical with some turns, while the last one is partially disordered and contains a kink in its geometry. We propose that these characteristics likely account for the lowest unfolding force observed, aligning with the predictions of the X-LoRA model.

In a second instance, we focus on the energetic aspects of proteins, specifically considering the unfolding energy required as the protein is extended from its initial folded state to a fully unfolded conformation~\cite{Ni2023ForceGen:Model}. We again analyze three sequences (the first two being identical to those mentioned earlier, while the third is a longer construct, combining the second sequence at the termini with the first sequence in the middle):

\begin{quote}
\small\texttt{\begin{itemize}
    \item LNTWFDKFDAEICIPYVEILKESGMK \item AITWFDKFDAEICIPYVEALKIAAMK \item AITWFDKFDAEICIPYVEALKIAAMKLNTWFDKFDAEICIPYVEILKESGMKAITWFDKFDAEICIPYVEALKIAAMK
\end{itemize}
}
\end{quote}

We initially request the model to calculate the unfolding energy for each sequence, followed by an analysis of the results: based on the amino acid sequence, we ask for an explanation of why this protein is expected to be the most stable. Finally, we inquire about the probable functions that the most stable protein might perform.

Figure~\ref{fig:MJB_20} presents a transcript of a dialogue that leverages combinations of different experts and synthesizes knowledge from these capabilities, demonstrated here in the context of unfolding energy analysis for three proteins. By the conclusion of the discussion, the CoT/reasoning adapter emerges as the dominant component, with the X-LoRA model responding to a prompt that involves reasoning over the results to predict probable protein structural features and potential protein functions. Specifically, the mechanistic inquiry regarding the basis of protein stability highlights the formation of hydrophobic interactions among nonpolar amino acids as well as hydrogen bonds between polar amino acids. The model forecasts that several hydrophobic amino acids (I, F, W, Y) and hydrogen-bonding amino acids (D, E, K) play a role in stabilizing the protein.

Additionally, the model anticipates that cysteine (C) residues present in the sequence will facilitate the creation of disulfide bonds, further reinforcing the protein’s structural stability. It is noteworthy that the crucial predictions related to structural characteristics explaining the high stability are accurately identified, as confirmed through structural analyses depicted in Fig.~\ref{fig:MJB_21}. This analysis validates the presence of all these features. It also clarifies why the other two proteins, which are considerably shorter segments with differing stability in their helical regions, exhibit substantially lower unfolding energies.

\subsection{Adversarial agentic modeling to connect distinct scholarly disciplines and knowledge yielding ontological knowledge graph generation}

Given that our X-LoRA model possesses extensive cross-domain expertise, it can be employed to investigate links among diverse concepts, knowledge repositories, and fields of specialization. In one experiment, we pose two queries to the model, each similarly formulated but emphasizing a different perspective. Subsequently, we generate triplets to construct an ontological knowledge graph, thereby transforming the responses into more structured representations~\cite{Giesa2012CategoryDesign,Spivak2011ReoccurringAnalogies}. The resulting graph delivers a cohesive understanding of the derived insights and visually maps connections between concepts in an interpretable and mechanistic way.

We employ two adversarial agents to address the question. Agent 1 is assigned the role of a "critical physicist," while Agent 2 acts as a "philosopher." Agent 1 initiates the inquiry and is responsible for probing and formulating additional questions in response to each answer given by Agent 2. As the dialogue progresses, the discussion delves deeper into the topic, revealing various facets. The question is deliberately selected to span multiple disciplinary domains, aiming to evaluate the model’s capacity to synthesize ideas and concepts.

The question is as follows:

\begin{quote}
\small\texttt{Consider an abstract representation of a classical musical composition as a graph of interacting notes, chords, melodies, rhythms and so on.

What happens if we would apply pressure, up to a point of failure, on this abstract concept of music?

Provide logical steps, mathematical reasoning, and discuss specifically the physics of the problem.}
\end{quote}

The complete conversation is provided in the Supplementary Material (see, SI Conv Example 3). For brevity, we present here only the summary and key conclusion, which encapsulate the main points:

\begin{quote}
\small\texttt{The discussion explores applying pressure to an abstract representation of a classical musical composition, treated as a physical system composed of interacting notes, chords, melodies, and rhythms. The pressure is measurable via stress, and the failure threshold is identified when stress reaches a critical magnitude. To verify the mathematical model’s predictions, a physical analogue of the abstract musical composition can be employed, allowing experimental results to be compared with theoretical expectations. The mathematical model’s limitations include a simplified portrayal of interactions, a static framework, an assumption of homogeneous distribution, and a linear correlation between stress and pressure. Enhancing the model’s fidelity requires incorporating more complex and realistic interactions, dynamic behaviors, heterogeneous distributions, and nonlinear terms.

[...]

The most significant insight is that applying pressure to this abstract representation of a classical musical composition can be conceptualized as altering the force field among the interacting notes, chords, melodies, and rhythms. Increasing pressure leads to heightened stress within the system, and beyond a critical threshold, the system fails and collapses under its own influence. This understanding addresses the original question by offering a physical rationale for how pressure impacts the abstract musical concept and how this effect can be quantified via stress.}
\end{quote}

Although the generated text offers valuable perspectives, constructing an ontological representation of the data can enhance interpretability~\cite{Buehler2023MechGPTModalities}. For graph construction, we combine the entirety of both conversations and generate a knowledge graph. An example of such a graph is displayed in Figure~\ref{fig:graph}. The resulting graph is partitioned into communities using the Girvan-Newman algorithm~\cite{Girvan2002CommunityNetworks}, yielding a total of 12 distinct communities. The figure presents both the complete graph (Figure~\ref{fig:graph}a) and detailed views with node annotations (Figure~\ref{fig:graph}b-c).

\subsection{Creation of X-LoRA-Gemma Integrating Protein, Chemical, Bio-inspired, and Materials Mechanics Functionalities}
\label{gemma-section}

To validate that our proposed methodology is applicable to different foundational models with varying architectures, we trained an additional X-LoRA model, this time utilizing the Gemma-7B-it model~\cite{Google/gemma-7b-itFace}. In this instance, we implemented a set of four LoRA adapters, defined as follows:

\begin{enumerate}
    \item Bioinspired materials \item Mechanics and materials \item Protein mechanics tasks (including generative sequence-to-property and inverse design capabilities) \item Quantum-mechanics based molecular properties from QM9 (incorporating generative SMILES-to-property and inverse design functionalities; refer to Table~\ref{tab:table_QM9} for definitions and summaries of all properties calculated within this dataset)~\cite{Ruddigkeit2012EnumerationGDB-17,Ramakrishnan2015ElectronicSpace}
\end{enumerate}

This X-LoRA-Gemma model was developed analogously to the original X-LoRA model but aims to incorporate a distinct range of functionalities, particularly adding the capacity to predict a suite of 12 quantum mechanical properties and to design molecules tailored to these 12 quantum mechanical attributes. This new adapter was trained alongside the other adapters mentioned above, which rely on the same datasets utilized in our previous experiments. It is important to note that while the main body of the paper concentrates on the Zephyr-based X-LoRA model, within this section, to differentiate between models, we refer to the original model as X-LoRA-Zephyr and the Gemma-based model as X-LoRA-Gemma.

Figure~\ref{fig:protein_force_energy_prediction} illustrates the overall performance of X-LoRA, demonstrating predictions of unfolding force and unfolding energy from sequence data ($R_2 = 0.87$), as well as predictions of quantum mechanical properties for molecules from the QM9 dataset ($R_2 = 0.96$). We observe excellent predictive accuracy on both tasks, with the protein design task performance being marginally better than that of the original X-LoRA model.

We next showcase the capabilities of X-LoRA-Gemma through a straightforward multi-agent design experiment. We examine a small organic molecule, \texttt{CCCC12CC1C(=O)N2} (drawn from the test set of the QM9 dataset). This molecule features a bicyclic framework with both a ketone (=O) and an amine (N) integrated into the ring structure; see Figure~\ref{fig:QM9_design}(a) for a depiction of the molecule. Initially, we pose the following prompt to the X-LoRA-Gemma model:

\begin{quote}
\small\texttt{Here is a molecule in SMILES code: CCCC12CC1C(=O)N2

What modifications would I need to make to this molecule to enhance its utility for polymer design?}
\end{quote}

The answer is:

\begin{quote}
\small\texttt{To enhance the utility of this molecule for polymer design, you might consider:
\begin{enumerate}
\item Adding a reactive functional group, like an amine or hydroxyl group, to facilitate further modification and crosslinking.
\item Extending the carbon backbone to make it longer and more flexible, which would improve the molecule's capability to form polymers with diverse structures and characteristics.
\item Introducing a side chain containing a functional group, such as a carboxyl group, to enable the formation of covalent bonds with other molecules and polymers.
\item Integrating a cyclic structure into the molecule, for example, a lactam ring, to generate a polymer that is more rigid and stable.
\item Incorporating various side chains to produce polymers with a broad range of properties and functionalities.
\item Varying the number of carbons in the backbone to facilitate the synthesis of polymers with differing molecular weights and attributes.
\item Adding a crosslinking site within the molecule to yield a polymer that is more stable and less reactive.
\item Designing a repeating unit structure to achieve a polymer with targeted molecular weight and properties.
\item Combining several of these approaches to develop a flexible and practical polymer suitable for multiple applications.
\end{enumerate}}
\end{quote}

We will now concentrate on the first recommendation.

\begin{quote}
\small\texttt{How can the overall reactivity of a molecule be increased? I intend to utilize this molecule for constructing a polymer material.

The molecule in question, expressed in SMILES notation, is: CCCC12CC1C(=O)N2

I have the ability to adjust the following characteristics:

Dipole moment, mu Polarizability, alpha Highest Occupied Molecular Orbital energy, homo Lowest Unoccupied Molecular Orbital energy, lumo HOMO-LUMO gap: The energy difference between HOMO and LUMO, gap Electronic spatial extent Zero-point vibrational energy Heat capacity at constant volume Internal energy at 0 K Internal energy at 298.15 K Enthalpy at 298.15 K Free energy at 298.15 K Which are the top THREE properties I should modify, and in what way (increase or decrease)? }
\end{quote}

The feedback from X-LoRA-Gemma is:

\begin{quote}
\small\texttt{\begin{itemize}
    \item  HOMO-LUMO gap (decrease) \item  Dipole moment (increase) \item  Polarizability (increase)
\end{itemize}
Adjusting these molecular characteristics can enhance the overall reactivity of the compound and render it more appropriate for use in polymer material synthesis. }
\end{quote}

We leverage the generative and predictive functionalities of X-LoRA-Gemma to design molecules exhibiting these targeted property modifications.

For the optimization, we implement a two-agent framework comprising two instances of X-LoRA-Gemma. One agent performs the generative task \texttt{GenerateMolecularProperties<...>} to produce candidate molecules, while the second agent executes the \texttt{CalculateMolecularProperties<...>} task to evaluate the discrepancy between the desired and the predicted properties. This iterative cycle continues until either the maximum iteration count is reached or the mean-squared error falls below a predefined threshold.

The outcomes of this procedure are illustrated in Figure~\ref{fig:QM9_design}. Panel (a) of Figure~\ref{fig:QM9_design} displays the target molecular design alongside its deviation from the original molecule's properties. Panel (b) presents the top 20 generated designs, accompanied by the mean-squared error distribution in panel (c) and a plot of the mean-squared error progression over the generated candidates in panel (d). The optimal design identified is \texttt{CC(C)C1=CC(=O)NC=C1}, which approximates the target closely, as shown in panel (e). Lastly, panel (f) depicts the optimized 3D molecular structure (optimized using UFF~\cite{Rappe1992UFFSimulations}).

To evaluate the molecule's reactivity, we also calculated Gasteiger charges~\cite{Gasteiger1980IterativeCharges} (see Figure~\ref{fig:QM9_design}(f) for a visual summary). The molecule contains an aromatic ring with a ketone (C=O) group, which is more electrophilic than typical carbon-carbon or carbon-hydrogen bonds, thereby rendering it prone to nucleophilic attack. The presence of both the nitrogen-containing aromatic ring and the ketone group creates a pyridone ring structure. This ring exhibits higher reactivity compared to a simple aromatic ring due to the heteroatom (nitrogen) and the ketone introducing reactive sites, such as for nucleophilic addition at the carbonyl carbon or electrophilic substitution on the ring. The nitrogen atom, as part of the heteroaromatic system, influences the aromaticity and reactivity by modifying the electron density distribution across the ring. This modulation impacts the molecule’s interactions with other chemical species, designating the nitrogen-containing ring as a site for potential reactions, especially those involving electrophilic attack facilitated by the electron-donating properties of the nitrogen.

The carbon atom within the ring, positioned next to both an oxygen and a nitrogen atom and carrying a partial positive charge (+0.25), indicates it is somewhat deficient in electrons. This arises because oxygen, being more electronegative, draws electron density toward itself, reducing the electron richness of the carbon. This influence is somewhat offset by the neighboring nitrogen, which is less electronegative than oxygen but more electronegative than carbon. Nonetheless, the partial positive charge signifies that the carbon is a probable target for nucleophilic attack, where an electron-rich nucleophile is drawn to this electron-poor carbon.

The nitrogen atom, exhibiting a partial negative charge (-0.33), suggests it is relatively electron-rich compared to its typical state. This is attributable to its lone pair of electrons and the impact of the aromatic ring’s electron system. In aromatic heterocycles, the nitrogen’s lone pair may contribute to the delocalized $\pi$-electron cloud, although the exact degree depends on the ring’s architecture and the electronegativities of neighboring atoms. The partial negative charge implies that nitrogen is more nucleophilic, possessing a greater tendency to donate electrons, either by forming bonds with electrophiles or by undergoing protonation. Additional quantum mechanical studies are necessary to explore this further.

The carbon carrying the partial positive charge represents a reactive center for nucleophilic attack. Various molecules could interact with this site via mechanisms involving nucleophiles targeting electron-deficient carbons.

The hydrogen atom bonded to the nitrogen, with a partial positive charge of 0.17, serves as a favorable site for hydrogen bonding with other atoms or molecules. Hydrogen bonds are dipole-dipole interactions occurring when a hydrogen atom bonded to a highly electronegative atom (such as nitrogen, oxygen, or fluorine) associates with another electronegative atom possessing a lone electron pair. This characteristic endows the molecule with the potential to act as a hydrogen bond donor, which plays a crucial role in determining the molecule’s solubility in water and other polar solvents, as well as its behavior in biological contexts.

The ketone group positioned next to the nitrogen-containing aromatic ring imparts distinctive reactivity that can be leveraged in polymer synthesis~\cite{McQuarrieSimonPhysicalChemistry,CareySundbergAdvancedOrganicChemistry,Varghese2022BeyondApplications,Varghese2022BeyondApplications}. This compound is classified as a lactam, due to the cyclic amide structure formed by the ketone and nitrogen within the ring. Lactams play a significant role in polymer chemistry, especially in producing polyamides, a category of polymers widely utilized in various material applications. The ketone functional group is capable of undergoing multiple chemical reactions relevant to polymer chemistry, such as nucleophilic addition, which can be harnessed to extend polymer chains or introduce side groups that alter polymer characteristics. The nitrogen atom embedded in the aromatic ring can markedly influence the polymer’s electronic properties and engage in hydrogen bonding, thereby affecting properties such as melting point, solubility, and mechanical behavior. N-heteroaromatic compounds are also recognized for their capability to coordinate with metal ions, potentially enabling uses in catalysis or materials with electronic or photonic functionalities. This molecule may also participate in polymerization through its amide (lactam) functionality; polyamides are known for their durability, elasticity, and resistance to wear and chemicals, and find applications across diverse areas including textiles (e.g., nylon) and automotive components~\cite{Varghese2022BeyondApplications,Varghese2022BeyondApplications}.

Alkyl substituents can enhance the hydrophobic nature of a molecule, suggesting that polymers synthesized from monomers containing such alkyl groups are likely to exhibit increased hydrophobicity. This influences their solubility in various solvents and their interaction with water, which can be beneficial for applications requiring water resistance or particular solubility traits. Additionally, the incorporation of alkyl groups may impact the polymer's mechanical characteristics. For example, they may function as plasticizers within the polymer matrix, thereby improving the polymer's flexibility. Depending on the alkyl groups' length and branching, this could result in materials that are more flexible or possess enhanced impact resistance.

\subsubsection{Additional benchmarking} We present further evaluations of the X-LoRA-Gemma model, illustrated in Figure~\ref{fig:perf_MMLU}, using a subset of the MMLU benchmark focused on chemistry, physics, and biology (domains closely related to those on which our model was trained)~\cite{Hendrycks2020MeasuringUnderstanding}. Our findings show that both X-LoRA models discussed herein outperform their respective base models. Specifically, the X-LoRA models surpass both base models (Zephyr model: 54.6\% for the base versus 56.6\% for the X-LoRA model; Gemma model: 40.6\% for the base and 52.4\% for the X-LoRA-Gemma model). The highest overall accuracy is achieved by the X-LoRA-Zephyr model, though both models exhibit relatively close performance.

\section{Conclusion}

There exist numerous current and potential future applications of multimodal LLMs in scientific fields, alongside a strong demand for highly specialized expertise tailored to specific purposes. Beyond general language processing tasks, the ability to perform computational or generative tasks, as exemplified here with protein mechanics, is crucial~\cite{Luu2023BioinspiredLLM:Materials,Buehler2023MechGPTModalities,Buehler2023MeLMProblemsb}. X-LoRA, which is based on open-source language models, allows for the dynamic integration of expertise and knowledge from various domains, as illustrated in multiple examples, including Figure~\ref{fig:conversation_evolve}. This investigation demonstrated how the X-LoRA model autonomously and dynamically adjusts the utilization of adapters and combines certain adapters in a sophisticated manner across the LoRA layers. This intricate blending of LoRA layers is consistently observed throughout all the examples examined here, such as in Figures~\ref{fig:QA_weights_sample_0} and \ref{fig:QA_weights}. This phenomenon also presents an intriguing challenge to deepen our understanding of the underlying mechanisms from a physics standpoint, where a LLM can be regarded as a large-scale graph-generating model that leverages complex graph structures. Consequently, the mixing of adaptation experts at deeper layers, as implemented here, modifies these mechanisms in ways that foster the development of effective solutions.

The introduced algorithm enhances the traditional inference method, which relies on a single forward pass, by incorporating two forward passes. This approach trains the model to initially ‘reflect’ on the question and consider how it might reconfigure itself before providing an answer. This represents a straightforward implementation of ‘self-awareness,’ enabling the model to adjust its own architecture to optimally address a task. As a result, despite having a relatively modest parameter size of 7 billion, the model can reason across a wide range of scientific disciplines (such as biological materials, mathematics, physics, chemistry, logic, mechanics, etc.), substantially boosting its ability to generate creative solutions. Moreover, this approach could inspire the design of alternative model architectures, for example, techniques that extend the reconfiguration strategy to portions or the entirety of a model that has not been adapted. In such cases, the scaling head would dynamically alter segments of a pretrained model or combine multiple models. As demonstrated by X-LoRA, such methods can serve as powerful tools to broaden or unify the capabilities of different models.

We conducted multiple performance evaluations, including a comparison with significantly larger models (Figure~\ref{fig:perf_comp_bioinspiredexam}), where X-LoRA demonstrated clear superiority. The comprehensive comparison presented in Table~\ref{tab:table_qa} revealed that X-LoRA significantly outperforms the base model across a wide range of tasks and application areas. Additional comparisons involved quantitative evaluations of numerical prediction tasks, such as protein mechanics and the quantum-mechanical properties of molecules, where X-LoRA achieved high accuracies with $R_2$ values ranging from 0.85 to 0.96. For context, these results surpassed previously published outcomes, such as the prediction of protein unfolding force using the MeLM model~\cite{Buehler2023MeLMProblemsb}, which reported an $R_2$ of 0.78. The ability to accurately address multiple prediction tasks, alongside inverse design tasks and advanced reasoning skills, was highlighted in the molecular design example (Section~\ref{gemma-section}). This example also demonstrated that X-LoRA can be effectively applied to different base model architectures (Gemma vs. Zephyr/Mistral). Future research could further extend this approach and develop X-LoRA models for larger base models.

A limitation of the proposed method is that it requires two forward passes, which may increase computational cost: one pass to compute hidden states used as inputs to the X-LoRA scaling head, and a second pass with $\Lambda$ applied to the adapters. However, the first pass is utilized for the X-LoRA scaling head, leveraging the fundamental capabilities of the LLM rather than learning knowledge separately by training a larger scaling head. For effective deployment in model-serving tools, distinct key-value caches should be maintained for both the scaling and forward passes. To enhance inference speed, we developed \texttt{Mistral.rs}, an LLM serving platform that implements X-LoRA in Rust~\cite{matsakis2014rust} and incorporates multiple performance optimization techniques. Nonetheless, a key advantage of our approach is its compatibility with any model without the need to restructure internal logic. We consider this beneficial, particularly because it can be applied to the extensive resources available within the Hugging Face ecosystem.

A key benefit of the proposed method is its straightforward implementation within any existing LLM (for example, since the X-LoRA code has been made available, it can be easily applied to any model within the Hugging Face ecosystem). Another benefit lies in the scaling head's ability to leverage the inherent capabilities of the LLM, enabling it to learn complex strategies for optimally combining adapters to generate specific answers. When training X-LoRA, it is ideal to use complex question-answer pairs or conversations that allow the model to discover the most effective combinatorial approach. We conducted various analyses related to protein design and mechanics. For example, Figure~\ref{fig:MJB_20} and Fig.~\ref{fig:MJB_21} showcase stability analyses of multiple protein sequences, including comparative assessments, reasoning over predicted outcomes, and mechanistic explanations that offer a bottom-up chemical basis for the predicted behaviors. As illustrated in Fig.~\ref{fig:MJB_21}, these findings were directly validated through structural analysis of the protein’s folded 3D configuration.

Regarding future directions, further research could investigate a broader range of adapter experts, expanding upon the foundational expertise present in our LoRA set. The X-LoRA model outperforms the base model or a base model equipped with only one adapter, as it can address complex questions by leveraging the specialized capabilities contributed by the ensemble of adapters. There also appear to be synergistic effects among the various adapters, as suggested by the intricate mixing patterns of scaling weights across layers. This phenomenon could be studied in greater depth and compared with techniques such as SLERP that integrate models using alternative strategies. While we trained the X-LoRA scaling head with a subset of the training data—consisting of several hundred samples from each of the original datasets used to develop the adapters—there is potential to create more purposefully curated training sets tailored to this objective. For instance, we observed that multiple-choice questions tend to trigger reasoning-focused adapters, given their training on such data. Additional prompting techniques may be necessary to guide the scaling head toward understanding specific domains relevant to the questions, possibly through particular system messages or other methods. Overall, designing appropriate training datasets is an important area for further study, including specialized methods to effectively encourage layer-wise scaling mixing. Our experiments tracking these mechanisms over extended conversations demonstrate that the model can dynamically adjust scaling strategies to optimally address certain tasks, which is an encouraging outcome.

We identified intriguing patterns in how experts are activated across different layers, with compelling insights elaborated upon in the paper. This analysis has the potential to be further developed and may provide significant understanding regarding the benefits of mixing components or sets of layers within models. Such exploration would be especially valuable when applying the X-LoRA technique to domains beyond protein mechanics and protein design, as demonstrated here. We defer this to future research.

Another key direction involves deeper exploration of agentic modeling, exemplified here by the interaction of two X-LoRA models in an adversarial setup. Future studies could involve more than two models to enrich interactions and introduce additional functionalities, thereby advancing models into novel generative territories. This approach may also facilitate the creation of methods that integrate physics-based modeling or incorporate other validation procedures by employing code-writing/executing agents or agents utilizing function calling and other processing mechanisms~\cite{Ghafarollahi2024ProtAgents:Learning}.

\section{Materials and Methods}

The X-LoRA codebase and example implementations can be accessed at \url{https://github.com/EricLBuehler/xlora}. Model weights and datasets relevant to the X-LoRA described in this paper, along with supplementary examples, are available at \url{https://huggingface.co/lamm-mit/x-lora}. The X-LoRA-Gemma model can be found at \url{https://huggingface.co/lamm-mit/x-lora-gemma-7b}. The \texttt{Mistral.rs} inference engine is hosted at \url{https://github.com/EricLBuehler/mistral.rs}.

\subsection{Mathematical formulation of X-LoRA} We define the $\mathbf{scalings}$ as a matrix comprising token- and layer-level coefficients for LoRA adapters. The X-LoRA $\mathbf{scaling}$ $\mathbf{head}$ receives the base model’s hidden states and processes them through linear fully-connected layers to predict scalings, followed by applying a softmax function along the final dimension to ensure that scalings across all experts sum to one. ReLU activation functions and dropout layers are interspersed between the linear fully-connected layers within the X-LoRA scaling head. The size of the hidden layers is configurable and can be adjusted either to reduce dimensionality from the hidden dimension to the scaling dimension or to implement a widening and narrowing scheme akin to the feed-forward layer within a transformer block.

The scalings are represented by a matrix $\Lambda \in \mathbb{R}^{s \times l \times n}$, where $s$, $l$, and $n$ correspond to the sequence length, number of layers, and number of adapters, respectively. For each layer, the relevant scalings are extracted as a matrix $\lambda \in \mathbb{R}^{s \times n}$. The scalings $\lambda$ are then applied to each LoRA adapter by broadcast multiplying the corresponding X-LoRA scaling $\lambda_i \in \mathbb{R}^{1\times 1 \times s}$ with the input to the decomposition matrices $A$ and $B$.

Given $W_0 \in \mathbb{R}^{d \times k}$, $B_i \in \mathbb{R}^{d \times r_i}$, and $A_i \in \mathbb{R}^{r_i \times d}$, where the rank $r_i \ll {\rm min} (d,k)$ and is specific to each adapter $i$, the forward pass of an X-LoRA layer is described by the equation:

\begin{equation}
    h = W_0x + \sum_{i=1}^{n} B_iA_i(x \cdot  \lambda_i) \alpha_i
\end{equation}

The main paper includes visualizations of the scalings, which were generated either through bar plots or by using the \texttt{contourf} function from Matplotlib~\cite{MatplotlibDocumentation}.

\subsection{X-LoRA implementation algorithm and code} The chosen implementation prioritizes user-friendliness and the flexibility to adapt easily to different models, especially those within the Hugging Face framework. The X-LoRA method functions by:

\begin{enumerate}
    \item Keeping the base model and adapters fixed while optimizing performance.
    \item Combining LoRA adapters based on the predictions generated by the X-LoRA scaling head.
\end{enumerate}

The entire codebase is developed using Python and PyTorch~\cite{Paszke2019PyTorchLibrary}.

\subsubsection{Dual forward pass strategy for self-aware inference} Since X-LoRA necessitates computing hidden states to predict scalings, we adopt a dual forward pass technique: the first pass allows the model to assess the task and determine how to adjust itself, followed by the actual inference step. This approach introduces a kind of ‘self-awareness,’ enabling X-LoRA to allocate extra computation time for reflection instead of immediate output.

We refer to the initial pass as the scaling pass and the subsequent one as the forward pass. During the scaling pass, we initialize the scalings $\Lambda$ to zero. We have also tested setting scalings $\Lambda$ to ${n}^{-1}$ and observed comparable outcomes, suggesting stable behavior (this parameter is configurable in the X-LoRA package). To facilitate the application of X-LoRA across various models, we incorporate a forward pass hook. This mechanism is critical in linking the forward and scaling passes and constitutes the core feature enabling compatibility with all models.

\begin{enumerate}
    \item $\mathbf{Scaling}$ $\mathbf{pass}$: Executed by running the base model with adapters fixed at constant values. The generated hidden states are then utilized by the X-LoRA scaling head to produce $\Lambda$.
    \item $\mathbf{Forward}$ $\mathbf{pass}$: During this pass, $\Lambda$ is applied to blend the adapters, and the resulting logits are output as the final predictions from the X-LoRA model.
\end{enumerate}

A cautionary note: using the same key-value cache for both the scaling pass and forward pass might disrupt the correct input processing for the scaling head because the key-value cache would persist and accumulate data across both passes. Therefore, we disable the key-value cache during training and inference, with appropriate management implemented in \texttt{Mistral.rs}.

\subsubsection{Freezing weights and selectively enabling parts of the model for training} The base model and adapters are kept frozen (i.e., their weights remain unchanged during training), making the X-LoRA scaling head the only component in an X-LoRA model that is trainable. The code provided also allows the possibility to unlock all adapters along with the X-LoRA scaling head for training. Additionally, it is feasible to train the entire system—including the X-LoRA scaling head, adapters, and the foundational base model—either simultaneously or with distinct learning rates. Furthermore, future developments of the concepts introduced here might explore scaling functions between multiple foundation models to realize a traceable implementation of SLERP-like modeling. This avenue will be addressed in subsequent research.

\subsubsection{X-LoRA layers} An X-LoRA layer functions by scaling and combining outputs from multiple adapters. It encapsulates the original LoRA layer and thus gains access to the specific adapters associated with that layer.

To seamlessly integrate X-LoRA into any model, the X-LoRA algorithm replaces the original LoRA layers by constructing X-LoRA layers, subsequently embedding their forward-pass logic into the LoRA layer. This approach leverages certain features of Python’s object-oriented model to facilitate swapping.

In brief, during the forward pass, each X-LoRA layer carries out the following steps (refer to Figure \ref{fig:general_arch}):

\begin{enumerate}
    \item $\mathbf{Scaling}$: Multiply the input to each LoRA adapter by its scaling factor $\lambda_i$.
    \item $\mathbf{Forward}$: Execute the forward propagation through the wrapped LoRA layer.
\end{enumerate}

\subsection{Optimized inference with \texttt{Mistral.rs} implemented in Rust} To accelerate inference, we created \texttt{Mistral.rs}, a Rust-based LLM serving platform that implements X-LoRA~\cite{matsakis2014rust} and incorporates multiple optimizations to boost performance. The following techniques are implemented, listed in no specific order:

\begin{itemize}
    \item LoRA adapter weight stacking: When all $A_i$ and $B_i$ matrices of each LoRA adapter share the same dimensions, the LoRA adapter weight matrices $A_i$ and $B_i$ can be stacked along their first dimension.

Mathematically, given $W_0 \in \mathbb{R}^{d \times k}$, for adapter matrices $A_i \in \mathbb{R}^{r \times d}$ and $B_i \in \mathbb{R}^{d \times r}$ with $r_i \ll {\rm min} (d,k)$, we stack the matrices so that $A\in \mathbb{R}^{i \times r \times d}$ and $B\in \mathbb{R}^{i \times d \times r}$. The $\alpha_i$ parameters are likewise applied during the initialization of these concatenated matrices.

To implement scalings, we rearrange the dimensions so that $\lambda \in \mathbb{R} ^ {i \times 1 \times 1}$. Afterwards, these scalings are applied via broadcast multiplication.

\item Non-granular scalings To minimize the number of executions of the scaling pass, we offer an option for non-granular scalings. This approach caches the scalings after $n$ completion steps, since the KV cache guarantees the sequence length will stay at one for the rest of the request. Utilizing this method greatly reduces latency by lowering the number of base model invocations.

\item Fused Compute Unified Device Architecture (CUDA) kernels To enhance the efficiency of the forward pass, we employ fused CUDA kernels. These fused kernels allow complex operations, such as Rotary Embedding (\cite{su2021roformer}), which would normally run sequentially, to be combined into a single kernel executed in parallel on a GPU. Our kernels are adapted from the vLLM implementation (\cite{kwon2023efficient}).

\item Quantization The inference engine supports quantization, enabling the use of a quantized base model. Preliminary experiments show that quantization to 8 bits performs well even with models originally trained using \texttt{bfloat16}, for example.

\end{itemize}

\subsection{Datasets}

Each adapter is trained on specialized datasets, summarized in Table~\ref{tab:table_loradata} along with references to their sources and original publications.

\subsection{Training strategy} Training proceeds in multiple stages. We start with the Zephyr-7B-$\beta$ model~\cite{HuggingFaceH4/zephyr-7b-betaFace}, which is based on the Mistral-7B model\cite{Jiang2023Mistral7B}, and train a series of adapters using the datasets introduced above. The Zephyr-7B-$\beta$ tokenizer is used throughout.

The initial eight adapters are trained using question-answer pairs provided by the respective datasets (see Table~\ref{tab:table_loradata}).

For the protein mechanics tasks, training follows a similar procedure but employs specific forward and inverse instruction sets. The bidirectional instruction tasks include:

\begin{quote}
\tiny {\texttt{CalculateForce< G E E C D C G S P ....> [0.310]\\
CalculateEnergy< G E E C D C G S P ....> [0.220]\\ 
CalculateForceHistory< G E E C D C G S P ....> [0.004,0.034,0.125,0.142,0.159,0.102,0.079,0.073,0.131, ...]\\
GenerateForce<0.310>\,[ G E E C D C G S P ....]\\
GenerateEnergy<0.220>\,[ G E E C D C G S P ....]\\
GenerateForceHistory<0.004,0.034,0.125,0.142,0.159,0.102,0.079,0.073,0.131, ...>\,[ G E E C D C G S P ....] }}
\end{quote}

These correspond to three distinct tasks in total: estimating the maximum unfolding force of a protein, calculating the unfolding energy, and determining the force-deformation trajectory. Each of these forward tasks is paired with a corresponding inverse task of the same type, which outputs a protein sequence instead, resulting in six tasks altogether.

The rank parameter for each LoRA adapter is set to $r=16$ with a scaling factor $\alpha=16$. Training of each adapter targets five specific modules: \texttt{v\_proj}, \texttt{k\_proj}, \texttt{q\_proj}, \texttt{gate\_proj}, and \texttt{o\_proj}, applying a dropout rate of 0.05. Typically, LoRA adapters are trained with relatively low ranks, ranging from 4 to 32, as it has been demonstrated that increasing the rank does not generally enhance performance but does raise computational demands. Lower ranks produce a more compact model with fewer parameters~\cite{hu2021lora}. The selected values for $r$, $\alpha$, and other hyperparameters are based on widely accepted defaults in the literature known to perform well. Additionally, prior research by the author has tested various configurations and found these particular choices to be effective~\cite{Luu2023BioinspiredLLM:Materials}.

The X-LoRA scaling head is trained on approximately 2,000 samples, which are a subset taken from the original dataset used for adapter training (this training set was also augmented with additional multiple-choice questions generated by GPT-4). We employ a \texttt{paged\_adamw\_8bit} optimizer with gradient clipping set to 0.3, a learning rate of $2\times 10^{-4}$ with a warmup phase, and four gradient accumulation steps, using a batch size of one. The X-LoRA model is trained for roughly 10,000 iterations, with the training loss settling around 0.28 at completion.

Given that the foundational Mistral architecture has 32 layers, and adapters are applied in 5 layers within each transformer layer, there are a total of 160 LoRA layers, each replaced by an X-LoRA layer. The X-LoRA scaling head generates scaling factors for each of these 160 LoRA layers (across nine LoRA experts, this yields $160 \times 9 = 1440$ $\lambda_i$ values computed per token). The scaling head consists of a single linear layer with ReLU activation and dropout set at $p=0.2$, totaling approximately 5 million parameters. The implementation is designed with flexibility to allow the development and evaluation of more complex deep classifier heads for other uses.

Both LoRA adapters and the X-LoRA scaling head are trained through next-token prediction using cross-entropy loss (we estimate the likelihood that the model assigns to the correct next token in a sequence, given the preceding tokens; the objective during training is to maximize this probability for the correct next token as indicated in the training dataset). We employ token shifting, where the sequence is shifted by one token to form the target sequence that the model should predict. For instance, if the input sequence is ``Proteins are fundamental to", the corresponding target sequence for training becomes ``are fundamental to biology". Thus, the model aims to predict the subsequent token in the sequence.

As illustrated in Figure~\ref{fig:hierarchical_concept}, during training of the X-LoRA scaling head, only the scaling head parameters are updated while all other parameters remain fixed. The native Zephyr tokenizer and prompt template are utilized.

\subsection{X-LoRA-Gemma model} The X-LoRA-Gemma model is developed in a manner similar to the X-LoRA model described above, but it is based on the Gemma-7B-it model~\cite{Google/gemma-7b-itFace} and incorporates four adapters (refer to the main text for details). These four adapters are trained on datasets related to mechanics and materials, protein mechanics, bioinspired materials, as well as an additional quantum-mechanics-based molecular properties QM9 dataset (see Table~\ref{tab:table_QM9})~\cite{Ruddigkeit2012EnumerationGDB-17,Ramakrishnan2015ElectronicSpace}.

Each LoRA adapter has a rank of $r=16$ with $\alpha=16$. Every adapter is trained targeting seven modules: \texttt{q\_proj}, \texttt{o\_proj}, \texttt{k\_proj}, \texttt{v\_proj}, \texttt{gate\_proj}, \texttt{up\_proj}, and \texttt{down\_proj}, applying dropout at 0.05. The rank for each adapter is consistently set at $r=16$ with $\alpha=16$. Given the four adapters and seven targeted modules, and with 28 layers in the Gemma-7B base model, a total of $28 \times 28 = 784$ values of $\lambda_i$ are computed per token.

The first two adapters (for bioinspired materials and mechanics of materials) incorporated into this model are trained using question-answer pairs as provided by their respective datasets (see Table~\ref{tab:table_loradata}).

For protein mechanics tasks, training is conducted similarly to the method described in the previous section, using both forward and inverse instruction sets in a bidirectional manner. Additionally, we train an adapter using the QM9 dataset with the following two tasks:

\begin{quote}
\tiny {\texttt{CalculateMolecularProperties<CC(C)N1CC1C=O> [0.098,0.358,0.581,0.395,0.309,0.330,0.570,0.649,0.519,0.519,0.519,0.518]\\
GenerateMolecularProperties<0.098,0.358,0.581,0.395,0.309,0.330,0.570,0.649,0.519,0.519,0.519,0.518> [CC(C)N1CC1C=O]\\ 
...
}}
\end{quote}

These correspond to two distinct tasks: calculating molecular properties (the forward task) and the inverse task of generating a molecule with specified target properties. Predictions or designs are made for all 12 properties included in the QM9 dataset (refer to Table~\ref{tab:table_QM9}).  
The model incorporates a total of four adapters.

The X-LoRA scaling head used in this model consists of three layers, with the intermediate layer size set to $3072 \times {FF}_{\rm fac} = 12,288$, where the feed-forward multiplier ${FF}_{\rm fac} = 4$. Linear layers with ReLU activation functions and a dropout rate of $p = 0.2$ are employed, resulting in approximately 47 million parameters.

The X-LoRA-Gemma scaling head is trained using roughly 17,000 samples. These samples were selected as a subset from the training dataset used for adapter training; notably, despite the X-LoRA-Gemma model containing only four adapters, the samples were drawn from the entire training set of the previously described model, supplemented with samples from the QM9 dataset. This approach was taken to explore whether the X-LoRA model might acquire additional capabilities during the scaling head training phase. Given that the loss converges well to 0.58 by the end of training, coupled with excellent performance even on tasks for which the model was not explicitly trained (such as chemistry, demonstrated in the molecule design example in the paper), we consider this approach to generally produce strong outcomes.

We employed a \texttt{paged\_adamw\_8bit} optimizer with gradient clipping set at 0.3, a learning rate of $1\times 10^{-4}$ accompanied by a warmup phase, and four steps of gradient accumulation, using a batch size of one. The X-LoRA-Gemma model was trained for approximately 62,000 steps.

The original Gemma tokenizer and prompt template were utilized.

\subsection{Adversarial agentic modeling} We implement an adversarial agentic approach by creating two X-LoRA agents. One agent specializes in asking questions: it poses the initial query and then continuously formulates further questions to uncover challenges and weaknesses in the responses. The other agent provides answers to these questions, consistently responding. This approach is implemented using \{Guidance\}~\cite{Guidance-ai/guidance:Models.}.

For the example connecting music and mechanics mentioned in the text, the roles of the two agents are defined as follows. First, the physicist and question asker:

\begin{quote}
\small\texttt{You are a critical physicist engaged in a discussion from the physics viewpoint. Keep your replies concise and consistently challenge statements in a provocative manner.

As a creative individual, you incorporate ideas from other disciplines and push boundaries.}
\end{quote}

Second, a philosopher (in some examples, this role is adjusted to represent a biology expert):

\begin{quote}
\small\texttt{You are a philosopher knowledgeable in biology, chemistry, and mathematics participating in a discussion.

Keep your answers concise, accurate, and creative. You generate excellent ideas and new directions of thought, always maintaining logical coherence.}
\end{quote}

The question asker is explicitly instructed to reflect on the current question and prior answers to develop new insightful queries. This is implemented via a prompt of the following form:

\begin{quote}
\small\texttt{Examine the following question and answer.

\#\#\# Question: <question>

\#\#\# Response: <response>

\#\#\# Instruction: Provide a SINGLE follow-up question that critically examines the response.

DO NOT provide an answer or commentary on the question at this stage.

The single question is: }
\end{quote}

After the dialogue has continued for $N$ exchanges, the model's tasks are to 1) summarize the main insights discussed, 2) present the key insights as bullet points, and 3) determine the most significant conclusion.

The raw conversation text serves as the foundation for subsequent analysis through graph creation.

\subsection{Knowledge graph generation} We employ Zephyr-7B-$\beta$ to extract triplets from the text, following the methodology described in~\cite{Rahulnyk/knowledge_graph:QnA}, enhanced with features inspired by the Llama Index graph generation method~\cite{Liu_LlamaIndex_2022}. The graphs are rendered using NetworX\cite{Networkx/networkx:Python} and Pyvis~\cite{WestHealth/pyvis:Graphs.}. One primary instruction, based on the approach in~\cite{Liu_LlamaIndex_2022}, is:

\begin{quote}
\small\texttt{Your objective is to extract the ontology of terms mentioned in the given context. These terms should capture the key concepts relevant to the context.

Present your output as a JSON list. Each list element should contain a pair of terms and their relationship, formatted like the following: [...] }
\end{quote}

We direct the construction of triplets comprising nodes and edge attributes as follows, adhering to the approach in\cite{Liu_LlamaIndex_2022}:

\begin{quote}
\small\texttt{\begin{itemize}
            \item "node\_1": "A concept from the extracted ontology"
            \item "node\_2": "A related concept from the extracted ontology"
            \item "edge": "A concise description of the relationship between node\_1 and node\_2"
        \end{itemize}
        }
\end{quote}

We also supply the model with examples of ontological analysis, such as:

\begin{quote}
\small\texttt{
       Context: ```Spider silk is a strong natural fiber used to catch prey in a web.``` }
\end{quote}

Sample triplets provided include:

\begin{quote}
\small\texttt{\begin{itemize}
            \item "node\_1": "spider silk"
            \item "node\_2": "fiber"
            \item "edge": "is"
        \end{itemize}
        }
\end{quote}

The assembled graph data is subsequently clustered into communities using the Girvan-Newman algorithm~\cite{Girvan2002CommunityNetworks}.

\subsection{Visualization of molecular structures}

We utilize PyMOL~\cite{Schrodinger2015The1.8} to visualize and examine the predicted protein structures. A variety of scripts and functions in this software assist in identifying protein characteristics such as secondary structure, hydrophobic and hydrophilic regions, disulfide bridges, and hydrogen bonds.

\section*{Data Availability Statement} The codes and datasets supporting the results of this study are freely accessible at \url{https://github.com/EricLBuehler/xlora} and \url{https://huggingface.co/lamm-mit/x-lora}. Additional data sources employed in this research are cited (see, specifically Table~\ref{tab:table_loradata}).

\end{document}